\section{Conclusions}
\label{sec:conclusions}

An edge-scattering problem has been posed and solved for a guided mode in a
$6mm$ n-type ZnO plate described by strain-gradient elasticity,
piezoelectricity, flexoelectricity and linearised electron drift--diffusion,
with one added carrier--strain-gradient enthalpy and its reciprocal partners.
Five results follow.

\begin{enumerate}
\item The energy identity closes at the reference plane to
$3.4\times10^{-16}$ relative, with the interference term quoted separately and
no evanescent member counted as transported power.
\item At $k_xd=0.736$ (\SI{16.68}{\giga\hertz}) the backward branch carries
\SI{2.31}{\percent} of the incident power; along the same branch the share rises
to a certified maximum of \SI{94.68}{\percent} at $k_xd=0.5275$
(\SI{16.6547}{\giga\hertz}), where the group velocity vanishes and the backward
wavenumber approaches the incident one to within \SI{0.4}{\percent}.  The
maximum is a local extremum of a single-humped, flat-topped response, stable to
$5.6\times10^{-9}$ between Galerkin degrees 14 and 34 and not claimed as a
global maximum.
\item These are modal shares and not reflectances: \SI{99.98}{\percent} of the
incident power leaves as propagating modes at the operating point and
\SI{99.77}{\percent} at the maximum, where the incident flux itself has fallen
by a factor $128$ because the branch is at a zero-group-velocity fold, so a
maximal share coexists with a minimal transported power.
\item The conversion is not produced by the added coupling.  With $H=0$ both
shares are unchanged, and removing the strain-gradient sector removes the
backward branch, leaving the specular reflection as the only negative-power
outgoing mode.
\item The coupling acts through conduction loss, reducing the partition by
\SI{1.21}{\percent} at the operating point and \SI{34.2}{\percent} at the
maximum, of which \SI{23.1}{} points are added propagation loss when the two
configurations are compared at the edge plane; re-searched at each coupling
level, the extremum shifts and its ceiling falls, so the coupling lowers the
achievable conversion as well as displacing it.
\end{enumerate}

The gradient length, the flexoelectric coefficient and the coupling tensor are
synthetic, the carriers are linearised about a uniform doping, the plate has a
single edge, corner line forces are not certified, and the relation to the
primary flexoelectric-semiconductor theory is established only against the
free-energy decomposition that could be inspected.  The observable that follows
is the modal partition of a slow guided branch near its zero-group-velocity
point.
